\section{Main Theorem}
\label{sec:main-theorem}

\begin{theorem}[Exact identity representation in the odd high-accuracy regime]
\label{thm:main}
Under Assumptions~\ref{assump:antipodal-oddness},
\ref{assump:high-accuracy}, and~\ref{assump:universal-sgd-success}, the
identity feature map exactly represents the entire class:
\[
\forall h\in\mathcal H\ \exists w_h\in\mathbb R^n\ \forall x\in\mathcal X,
\qquad
\operatorname{sign}_{s_0}(\langle w_h,x\rangle)=h(x).
\]
Moreover, the law
$\mathcal P_{\mathrm{id}}:=\delta_{\varphi_{\mathrm{id}}}$ is fixed before
every distribution and target and satisfies
\[
\forall\mathcal D\in\Delta(\mathcal X)\ \forall h\in\mathcal H,
\quad
\Pr_{\varphi\sim\mathcal P_{\mathrm{id}}}\!\left[
\inf_{w\in\mathbb R^n}
\Pr_{x\sim\mathcal D}\!\left[
\operatorname{sign}_{s_0}(\langle w,\varphi(x)\rangle)h(x)<0
\right]=0\right]=1.
\]
Consequently,
\[
\operatorname{dc}^{1/2}(\mathcal H)
\leq\operatorname{dc}(\mathcal H)\leq n\leq S\leq TS.
\]
These conclusions also hold when $\mathcal H=\varnothing$, with the target
quantifiers interpreted vacuously.  The result is deterministic after the
source premise, nonasymptotic, and uses the fixed finite horizon $T$.  The
displayed constants $1$ and $2$ have no hidden dependence on
$n,m,S,T,\eta,\varepsilon,\mathcal H,\mathcal D,h$, initialization, or the
sample path.
\end{theorem}

The theorem is a depth-two, bias-free, antipodally odd, strict-accuracy
specialization.  It does not assert the corresponding unrestricted-depth,
non-odd, or unrestricted-accuracy statement.
